\section{Related Work}
\label{sec:related}

\paragraph{PII detection benchmarks.}
PII-Bench~\cite{shen2025piibench} evaluates query-aware masking over 55 English PII types and is the only prior work with an explicit multi-subject split, on which LLMs degrade. REDACT~\cite{vats2026redact} stratifies 25 languages by sensitivity tier and disclosure form; Mind the Gap~\cite{zafar2026mindthegap} stress-tests English detectors under seven distribution shifts. Off-the-shelf filters degrade sharply on non-Latin scripts~\cite{uppala2026opf}. Our benchmark instead combines Korean financial categories with explicit input transformations and a jurisdiction-specific annotation policy.

\paragraph{De-identification.}
The Text Anonymization Benchmark~\cite{pilan2022tab} introduced direct/quasi-identifier distinctions and risk-weighted recall, motivating separate analysis of legal identifiers and contextual attributes. LLM-based anonymisers~\cite{staab2025anonymizers,yang2025rupta} and re-identification benchmarks~\cite{krco2026ratbench} evaluate privacy--utility trade-offs in English; Pasch and Cha~\cite{pasch2025privacyutility} study pseudonymisation with back-mapping, a distinct task from the all-mentions extraction recall measured here. The closest financial de-identification work~\cite{albanese2026anonymous} is English and industry-authored.

\paragraph{Korean financial language evaluation.}
TWICE~\cite{hwang2025twice} introduces KorFinMTEB and shows that translated English benchmarks can miss Korean financial semantics. NMIXX~\cite{lee2025nmixx} studies domain-adapted Korean--English embeddings and introduces KorFinSTS for financial semantic similarity. KFinEval-Pilot~\cite{hwang2025kfineval} evaluates Korean financial knowledge, legal reasoning and toxicity; KRX Bench~\cite{son2024krx} constructs company-knowledge questions covering Korean and other markets. These studies motivate domain- and language-specific evaluation. \kiii{} complements their semantic and knowledge tasks with typed PII spans, statutory category mappings and controlled input variation.

\paragraph{Korean PII.}
KDPII~\cite{fei2024kdpii} and Jang et al.~\cite{jang2024koreanpiiner} annotate synthetic Korean dialogues with a flat 33-tag scheme and report that models recognise universal PII (email, phone) far better than Korean-specific formats. Thunder-DeID~\cite{hahm2025thunderdeid} builds a 729-label taxonomy from 6{,}700 court judgments; its financial identifiers (accounts, cards, bills) are leaves of a generic ``unique number'' bucket without format modelling, and monetary, transaction and credit attributes are absent because judgments do not contain them. Our taxonomy inverts that emphasis and excludes public institution and product names.

\paragraph{Contextual privacy and misuse.}
Contextual-integrity benchmarks~\cite{mireshghallah2024confaide,shao2024privacylens,fan2024goldcoin,li2025privacibench} ground privacy judgments in norms or in HIPAA/GDPR; attribute-inference attacks~\cite{staab2024beyond} and automated spear-phishing studies~\cite{heiding2024spearphishing} motivate our Identifiability tier but evaluate neither Korean nor finance.
